\topic{本章实验、故意破坏与练习}

\begin{experimentbox}{从路径算到磁盘块}
选择一个小文件，记录目录项、inode 号、inode block、数据块号和块内偏移。
用镜像查看工具核对每个小端字段。
\end{experimentbox}

\begin{faultinjectionbox}{位图标空闲但 inode 仍引用}
在镜像副本中清除一个仍被文件引用的数据块位。预期 fsck 报告交叉所有权或
引用空闲块，正常挂载不得把该块再次分配。
\end{faultinjectionbox}

\begin{pitfallbox}
磁盘 struct 不能依赖 C++ padding；目录项名字需要长度边界；块号和 byte 偏移
要乘块大小转换。
\end{pitfallbox}

\textbf{练习}
\begin{enumerate}[leftmargin=2.2em]
  \item 4 KiB 块的 block 12 从哪个 byte 偏移开始？
  \item inode 号是否等于数据块号？
  \item 为什么删除目录项后还要更新 inode/位图？
\end{enumerate}
\begin{answerbox}
1. \texttt{0xC000}；2. 不等于；3. 名字、对象和空间分配是不同盘面结构，
必须一致更新。
\end{answerbox}
